% GENERATED FILE. Edit canonical lessons, then run npm run generate:books.
% canonical-source-hash: fnv1a64:1f3e0323674cd5fb
% canonical-lessons: FR-C488-calin, FR-C488-caliner, FR-C488-complicite, FR-C488-chuchoter, FR-C488-confier

\chapter{Câlins et confidences}
\label{ch:fr-calins-et-confidences}
\hlchaptermodality{488}

\begin{chapteropening}
\textbf{By the end of this chapter:} \emph{I can talk about closeness and sharing secrets.}

Complete the last lesson of chapter 488: I can talk about closeness and sharing secrets.
\end{chapteropening}

\section[le câlin]{le câlin — cuddle, hug}

\label{lesson:FR-C488-calin}



\begin{quote}
Before the new one: say the French for tenderness, then the French for to stroke.
\end{quote}

\subsection*{You'll want to know: le câlin}
\textbf{le câlin} — \textquotedblleft{}cuddle, hug\textquotedblright{}.

\begin{quote}
Avant de dormir, la petite veut un câlin.
\end{quote}

\begin{quote}
\emph{Before sleeping, the little girl wants a cuddle.}
\end{quote}

\emph{Faire un câlin à quelqu'un} — to give someone a cuddle. In Québec also \emph{un colleux}.

\begin{cousinweb}
\emph{le câlin} is built on \emph{câliner}. This is usually traced to a western dialect verb for resting in the shade or the warmth, which goes back to Latin \emph{calēre} \textquotedblleft{}to be warm\textquotedblright{} (Lewis \& Short); the origin is not fully settled.
\end{cousinweb}

\begin{grammarlens}[title={What you've built}]
One more word for telling, explaining and describing at B1.
\end{grammarlens}

\subsection*{Your turn}
\begin{itemize}
\raggedright
  \item \emph{Say it:} \emph{le câlin}
  \item \emph{Say it:} \emph{le câlin}, once more
  \item \emph{Say it:} a sentence of your own with \emph{le câlin}
  \item \emph{Recall:} say the French for tenderness, then the French for to stroke, then say \emph{le câlin} again
\end{itemize}

\begin{tcolorbox}[breakable,colback=teal!4,colframe=teal!35!black,title={Before you move on}]
What does le câlin mean? (\textquotedblleft{}Cuddle, hug\textquotedblright{}.) Say it once more.
\end{tcolorbox}
\section[câliner]{câliner — to cuddle}

\label{lesson:FR-C488-caliner}



\begin{quote}
Before the new one: say the French for to stroke, then the French for cuddle.
\end{quote}

\subsection*{You'll want to know: câliner}
\textbf{câliner} — \textquotedblleft{}to cuddle\textquotedblright{}.

\begin{quote}
Le père câline son bébé sur le canapé.
\end{quote}

\begin{quote}
\emph{The father cuddles his baby on the sofa.}
\end{quote}

\emph{Câlin} as an adjective means affectionate: \emph{un chat câlin}.

\begin{cousinweb}
\emph{câliner} is usually traced to a western dialect verb for resting in the shade or the warmth, which goes back to Latin \emph{calēre} \textquotedblleft{}to be warm\textquotedblright{} (Lewis \& Short); the origin is not fully settled.
\end{cousinweb}

\begin{grammarlens}[title={What you've built}]
One more word for telling, explaining and describing at B1.
\end{grammarlens}

\subsection*{Your turn}
\begin{itemize}
\raggedright
  \item \emph{Say it:} \emph{câliner}
  \item \emph{Say it:} \emph{câliner}, once more
  \item \emph{Say it:} a sentence of your own with \emph{câliner}
  \item \emph{Recall:} say the French for to stroke, then the French for cuddle, then say \emph{câliner} again
\end{itemize}

\begin{tcolorbox}[breakable,colback=teal!4,colframe=teal!35!black,title={Before you move on}]
What does câliner mean? (\textquotedblleft{}To cuddle\textquotedblright{}.) Say it once more.
\end{tcolorbox}
\section[la complicité]{la complicité — closeness, bond; complicity}

\label{lesson:FR-C488-complicite}



\begin{quote}
Before the new one: say the French for cuddle, then the French for to cuddle.
\end{quote}

\subsection*{You'll want to know: la complicité}
\textbf{la complicité} — \textquotedblleft{}closeness, bond; complicity\textquotedblright{}.

\begin{quote}
Il y a une grande complicité entre les deux sœurs.
\end{quote}

\begin{quote}
\emph{There is a strong bond between the two sisters.}
\end{quote}

Usually warm: a close understanding. In law, complicity in a crime.

\begin{cousinweb}
\emph{la complicité} is built on \emph{complice} \textquotedblleft{}partner, accomplice\textquotedblright{} (EtymDB), from Late Latin \emph{complex}, \emph{complicis} \textquotedblleft{}closely joined\textquotedblright{}, built on \emph{plicāre} \textquotedblleft{}to fold\textquotedblright{} (Lewis \& Short).
\end{cousinweb}

\begin{grammarlens}[title={What you've built}]
One more word for telling, explaining and describing at B1.
\end{grammarlens}

\subsection*{Your turn}
\begin{itemize}
\raggedright
  \item \emph{Say it:} \emph{la complicité}
  \item \emph{Say it:} \emph{la complicité}, once more
  \item \emph{Say it:} a sentence of your own with \emph{la complicité}
  \item \emph{Recall:} say the French for cuddle, then the French for to cuddle, then say \emph{la complicité} again
\end{itemize}

\begin{tcolorbox}[breakable,colback=teal!4,colframe=teal!35!black,title={Before you move on}]
What does la complicité mean? (\textquotedblleft{}Closeness, bond; complicity\textquotedblright{}.) Say it once more.
\end{tcolorbox}
\section[chuchoter]{chuchoter — to whisper}

\label{lesson:FR-C488-chuchoter}



\begin{quote}
Before the new one: say the French for to cuddle, then the French for closeness.
\end{quote}

\subsection*{You'll want to know: chuchoter}
\textbf{chuchoter} — \textquotedblleft{}to whisper\textquotedblright{}.

\begin{quote}
Les deux amies chuchotent au fond de la classe.
\end{quote}

\begin{quote}
\emph{The two friends are whispering at the back of the classroom.}
\end{quote}

\emph{Un chuchotement} — a whisper. \emph{Murmurer} is close.

\begin{cousinweb}
\emph{chuchoter} is usually described as an imitative word: it copies the soft sound of whispering, and the origin is not settled beyond that.
\end{cousinweb}

\begin{grammarlens}[title={What you've built}]
One more word for telling, explaining and describing at B1.
\end{grammarlens}

\subsection*{Your turn}
\begin{itemize}
\raggedright
  \item \emph{Say it:} \emph{chuchoter}
  \item \emph{Say it:} \emph{chuchoter}, once more
  \item \emph{Say it:} a sentence of your own with \emph{chuchoter}
  \item \emph{Recall:} say the French for to cuddle, then the French for closeness, then say \emph{chuchoter} again
\end{itemize}

\begin{tcolorbox}[breakable,colback=teal!4,colframe=teal!35!black,title={Before you move on}]
What does chuchoter mean? (\textquotedblleft{}To whisper\textquotedblright{}.) Say it once more.
\end{tcolorbox}
\section[confier]{confier — to entrust; (se confier) to confide}

\label{lesson:FR-C488-confier}



\begin{quote}
Before the new one: say the French for closeness, then the French for to whisper.
\end{quote}

\subsection*{You'll want to know: confier}
\textbf{confier} — \textquotedblleft{}to entrust; (se confier) to confide\textquotedblright{}.

\begin{quote}
Je me confie toujours à ma meilleure amie.
\end{quote}

\begin{quote}
\emph{I always confide in my best friend.}
\end{quote}

\emph{Confier un secret à quelqu'un}; \emph{se confier à quelqu'un}.

\begin{cousinweb}
\emph{confier} goes back to Latin \emph{cōnfīdere} \textquotedblleft{}to trust fully\textquotedblright{}, reshaped after \emph{fier}; both are built on \emph{fīdere} \textquotedblleft{}to trust\textquotedblright{} (Lewis \& Short).
\end{cousinweb}

\begin{grammarlens}[title={What you've built}]
One more word for telling, explaining and describing at B1.
\end{grammarlens}

\subsection*{Your turn}
\begin{itemize}
\raggedright
  \item \emph{Say it:} \emph{confier}
  \item \emph{Say it:} \emph{confier}, once more
  \item \emph{Say it:} a sentence of your own with \emph{confier}
  \item \emph{Recall:} say the French for closeness, then the French for to whisper, then say \emph{confier} again
\end{itemize}

\begin{tcolorbox}[breakable,colback=teal!4,colframe=teal!35!black,title={Before you move on}]
What does confier mean? (\textquotedblleft{}To entrust; (se confier) to confide\textquotedblright{}.) Say it once more.
\end{tcolorbox}
